\documentclass[10pt]{article}

\usepackage[utf8]{inputenc}

\usepackage[T1]{fontenc}

\usepackage{amsmath}

\usepackage{amsfonts}

\usepackage{amssymb}

\usepackage[version=4]{mhchem}

\usepackage{stmaryrd}

\usepackage{bbold}

\usepackage{graphicx}

\usepackage[export]{adjustbox}

\graphicspath{ {./images/} }



\begin{document}

р и цы $\quad\left\|a_{i k}\right\|_{1}^{n}-$ такой матрицы, что $\sum_{i, k=1}^{n} a_{i k} x_{i} \bar{x}_{k}$ есть эрмитова П. о. ф.; 2) положительно определенного ядра - такой функции $K(x, y)=K(y, x)$, что



\[

\int_{-\infty}^{\infty} \int_{-\infty}^{\infty} K(x, y) \varphi(x) \overline{\varphi(y)} d x d y \geqslant 0

\]



для любой функции $\varphi(x)$ с интегрируемым квадратом; 3) положительно определенной функции - такой функции $f(x)$, что ядро $K(x, y)=f(x-y)$ является положительно определенным. Класс непрерывных положительно определенных функций $f(x)$ с $f(0)=1$ совпадает с классом характеристических функций законов распределения случайных величин. БCЭ-3.



ПОЛОЖИТЕЛЬНО ОПРЕДЕЛЕННАЯ ФУНКЦИЯ комплекснозначная функция $\varphi$ на группе $G$, удовлетворяющая неравенству



\[

\sum_{i, j=1}^{m} \alpha_{i} \overline{\alpha_{j}} \varphi\left(x_{j}^{-1} x_{i}\right) \geqslant 0

\]



для любых наборов $x_{1}, \ldots, x_{m} \in G, \alpha_{1}, \ldots, \alpha_{m} \in \mathbb{C}$. Совокупность П. о. ф. на $G$ образует конус в пространстве $M(G)$ всех ограниченных функций на $G$, замкнутый относительно операций умножения и комплексного сопряжения.



Причина выделения этого класса функций состоит в том, что именно П. о. ф. определяют положительные бункционалы на групповой алгебре $\mathbb{C} G$ и унитарные представления группы $G$. Точнее, пусть $\varphi: G \rightarrow \mathbb{C}-$ произвольная функция и $l_{\varphi}: \mathbb{C} G \rightarrow \mathbb{C}-$ функционал, заданный равенством



\[

l_{\varphi}\left(\sum_{g \in G} \alpha_{g} g\right)=\sum_{g \in G} \varphi(g) \alpha_{g} g,

\]



тогда для положительности $l_{\oplus}$ необходимо и достаточно, чтобы $\varphi$ была П. о. ф. Далее, $l_{\varphi}$ определяет *-представление алгебры $\mathbb{C} G$ в гильбертовом пространстве $H_{\varphi}$ и, следовательно, унитарное представление $\pi_{\varphi}$ группы $G$, причем $\varphi(g)=\left(\pi_{\varphi}(g) \xi, \xi\right)$ для нек-рого $\xi \in H_{\varphi}$. Обратно, для любого представления $\pi$ и любого вектора $\xi \in H_{\varphi}$ функция $g \rightarrow(\pi(g) \xi, \xi)$ является П. о. ф.



Если $G$ - топологич. группа, то представление $\pi_{\varphi}$ слабо непрерывно тогда и только тогда, когда П. о. $ф$. непрерывна. Если $G$ локально компактна, то непрерывные П. о. ф. взаимно однозначно соответствуют положительным функционалам на $L^{1}(G)$.



Для коммутативных локально компактных групп класс П. о. ф. совпадает с классом преобразований Фурье конечных положительных мер на двойственных группах. Имеется аналог этого утверждения для компактных групп: непрерывная функция $\varphi$ на компактной группе $G$ является П. о. ф. тогда и только тогда, когда ее преобразование Фурье $\hat{f}(b)$ принимает положительные (операторные) значения на каждом әлементе двойственного объекта, т. е.



\[

\int_{G} f(g)(\sigma(g) \xi, \xi) d g \geqslant 0

\]



для всякого представления $\sigma$ и всякого вектора $\xi \in \pi_{\sigma}$.



Лuт.: [1] Х ь юи т т Э., Р о с с К., Абстрактный гармонический анализ, пер. с англ., т. 2 , М., 1975; [2] Н а й м а р к М. А., Нормированные кольца, 2 изд., М., 1968.



ПОЛОЖИТЕЛЬНО ОПРЕДЕЛЕННОЕ С. ПДРО комплекснозначная функция $K$ на $X \times X$, где $X-$ произвольное множество, удовлетворяющая условию



\[

\sum K\left(x_{i}, x_{j}\right) \lambda_{i} \bar{\lambda}_{j} \geqslant 0

\]



для любых $i, j \in \mathbb{Z}, \lambda_{i} \in \mathbb{C}, x_{i} \in X$. Измеримые П. о. я. на пространстве с мерой $(X, \mu)$ соответствуют положительным интегральным операторам в $L^{2}(X, \mu) ;$ включение в схему такого соответствия произвольных по- ложительных операторов требует введения обобшенных П. о. я., ассоциированных с оснащенными гильбертовыми пространствами (см. [1]).



Теория П.о.я. обобщает теорию положительно определенных бункций на группе: для положительной определенности функции $f$ на группе $G$ необходимо и достаточно, чтобы функция $K(x, y)=f\left(x y^{-1}\right)$ на $G \times G$ была П.о. я. В частности, на П.о.я. распространяются нек-рые результаты теории положительно определенных функций. Напр., т е о р е м а Б о х н еp a o том, что всякая положительно определенная функция есть преобразование Фурье положительной меры (т. е. интегральная линейная комбинация характеров), обобщается следующим образом: всякое П. о. я. (обобщенное) допускает интегральное представление с помощью т. н. әлементарных П. о. Я. относительно данного диффференциального выражения [1].



Лum.: [1] Б е p е з а н с к и й Ю. М., Разложение по соб-



\includegraphics[max width=\textwidth, center]{2023_07_08_d05ad46427b77c411f1bg-1}

ІОЛОЖИТЕЛЬНО ОПРЕДЕЛЕННЫИ ОПЕРАТОР - симметричный оператор $A$ в гильбертовом пространстве $H$ такой, что



\[

\left.\inf \frac{\langle A x, x\rangle}{\langle x, x\rangle}\right\rangle 0

\]



для любого $x \in H, x \neq 0$. Всякий П. о. с. является положительным оператором.



ПОЛОЖИТЕЛЬНОЕ РАССЛОЕНИЕ - обобщение понятия дивизора положительной степени на римановой поверхности. Голоморфное векторное расслоение $E$ над комплексным пространством $X$ наз. п о л о ж пт е л в н ы м (обозначается $E>0$ ), если в $E$ существует такая әрмитова метрика $h$, что функция



$v \longmapsto-h(v, v)$



на $E$ строго псевдовыпукла вне нулевого сечения. Если $X$ - многообразие, то условие положительности выражается в терминах кривизны метрики $h$. А именно, кривизны борже метрики $h$ в расслоении $E$ отвечает эрмитова квадратичная форма $\Omega$ на $X$ со значениями в расслоении Herm $E$ эрмитовых эндоморфизмов расслоения $E$. Условие положительности эквивалентно тому, что $\Omega_{x}(u)$ - положительно определенный оператор в $E_{x}$ для любого $x \in X$ и любого ненулевого $u \in T_{X}, x$.



В случае, когда $E$ - расслоение на комплексные прямые над многообразием $\dot{X}$, условие положительности равносильно положительной определенности матрицы



\[

\left\|-\frac{\partial^{2} \log h}{\partial z_{\alpha} \overline{\partial z}_{\beta}}\right\|,

\]



где $z_{1}, \ldots, z_{n}$ - локальные координаты на $X, h>0$ функция, задающая эрмитову метрику при локальной тривиализации расслоения. Если $X$ компактно, то расслоение на комплексные прямые $E$ над $X$ положительно тогда и только тогда, когда Чжэня класе $c_{1}(E)$ содержит замкнутую форму вида



\[

i \sum_{\alpha, \beta} \varphi \alpha \beta d z_{\alpha} \wedge d \bar{z}_{\beta} \text {, }

\]



где $\|\varphi \alpha \beta\|$ - положительно определенная әрмитова матрица. В частности, если $X-$ риманова поверхность, то расслоение над $X$, определяемое дивизором степени $d$, положительно тогда и только тогда, когда $d>0$. В случае, когда $E$ - расслоение ранга $>1$ над многообразием $X$ размерности $>1$, рассматривается также следующий более узкий класс П. р.: расслоение $E_{\rightarrow} X$ наз. п о ло о и т е ль ны м в смы сле Н ак а п о, если на $E$ существует такая эрмитова метрика $h$, что эрмитова квадратичная форма $H$ на расслоении





\end{document}